\section{Related Work}
\label{sec:related}

\subsection{Dynamic and Adaptive Inference}
\label{sec:related_dynamic}
Adaptive inference reduces redundant computation by spending compute as a function
of the input. Early-exit networks such as
MSDNet~\cite{huang2018msdnet}, RANet~\cite{yang2020ranet}, and
EENet~\cite{ilhan2024eenet} attach intermediate classifiers and terminate once a
confidence criterion is met. A complementary line equips a single backbone with
multiple execution widths or depths, such as slimmable~\cite{yu2019slimmable} and
Once-for-All~\cite{cai2020ofa} networks, and layer- or block-skipping methods that
drop residual units at inference~\cite{wang2018skipnet,wu2018blockdrop}. These
approaches share a common goal: deciding \emph{how much} compute to spend per
input. Our work adopts the same dynamic-depth backbone (a shared base/super pair)
but focuses on the orthogonal question of \emph{how the routing decision itself is
learned}.


\subsection{Learned Routing with Efficiency Objectives}
\label{sec:related_routing}
A large body of work makes the routing decision \emph{learnable} and trains it
jointly with an efficiency penalty. Gumbel-Softmax--based
routers such as AR-Net~\cite{meng2020ar} sample discrete execution choices
through a continuous relaxation~\cite{jang2017gumbel}; gating and
cosine-similarity controllers, NAS-style methods with cost-aware terms (e.g.,
DA-NAS~\cite{dai2020danas}), and channel-selection approaches with
activation/sparsity losses~\cite{herrmann2020channel} follow the same template.
These methods are almost universally trained with a combined objective
$\lambda_{\mathrm{acc}}L_{\mathrm{acc}} + \lambda_{\mathrm{flops}}L_{\mathrm{flops}}
(+\,\lambda_{\mathrm{uni}}L_{\mathrm{uni}})$, which has three documented
limitations.
\emph{(1) Objective conflict}: the accuracy and efficiency terms pull in opposite
directions, and their balance is fragile.
\emph{(2) Collapse}: a discrete, sampled policy tends to converge to a single
action, requiring an entropy or uniformity term (and additional hyperparameters)
to keep both paths alive.
\emph{(3) Single operating point}: the weight $\lambda$ fixes one
accuracy--efficiency trade-off in the trained weights, so serving a different
compute budget requires retraining.
Our regression formulation removes the efficiency weight and the anti-collapse
term entirely: because the decision is recovered by thresholding a regressed
scalar at deployment, a single model spans the whole trade-off family. 



\subsection{Efficient Object Detection}
\label{sec:related_detection}
Efficient object detection has been pursued through lightweight backbones,
quantization, and pruning, but these yield a static accuracy--efficiency point. 
We instead build on an AnyDepth-YOLO~\cite{kang2026anydepth} whose backbone
can be executed at two depths sharing the same weights, exposing a per-input depth
choice at no extra storage cost. 
Detection is a particularly favorable setting for per-frame routing: in driving video the scene difficulty varies sharply over time (dense urban intersections versus empty highways, day versus night), so the value of the deep path is highly input-dependent~\cite{pixels2027}, and a router can recover most of the deep path's accuracy while running the shallow path on the majority of frames. 
This motivates evaluating routing not on still images but on tracking video, as we do on KITTI~\cite{geiger2012kitti} and BDD100K~\cite{yu2020bdd}.